% GENERATED FILE. Edit canonical lessons, then run npm run generate:books.
% canonical-source-hash: fnv1a64:f7ee9eff60849585
% canonical-lessons: GU-C83-fifty, GU-C83-hundred, GU-C83-thousand, GU-C83-howmuch, GU-C83-which

\chapter{Fifty to a Thousand, and How Much}
\label{ch:gu-fifty-to-a-thousand-and-how-much}
\hlchaptermodality{83}

\begin{chapteropening}
\textbf{By the end of this chapter:} \emph{I can say fifty, a hundred and a thousand, and ask how much and which one.}

Complete the last lesson of chapter 83: I can say fifty, a hundred and a thousand, and ask how much and which one.
\end{chapteropening}

\section[\texorpdfstring{\gu{પચાસ} (pacās)}{પચાસ (pacās)}]{\gu{પચાસ} (pacās) — fifty}

\label{lesson:GU-C83-fifty}



\begin{quote}
Before the new one: say the Gujarati for thirty, then the Gujarati for forty.
\end{quote}

\subsection*{You'll want to know: \gu{પચાસ}}
\textbf{\gu{પચાસ}} (\emph{pacās}) — \textquotedblleft{}fifty\textquotedblright{}.

\begin{cousinweb}
\textbf{\gu{પચાસ}} (\emph{pacās}) is a tadbhava, a word inherited from Sanskrit and reshaped by sound change. It comes from \emph{pañcāśat}, \textquotedblleft{}fifty\textquotedblright{} (CDIAL 7682).
\end{cousinweb}

\begin{grammarlens}[title={What you've built}]
One more word for fifty to a thousand, and how much.
\end{grammarlens}

\subsection*{Your turn}
\begin{itemize}
\raggedright
  \item \emph{Say it:} \emph{pacās}
  \item \emph{Say it:} \emph{pacās}, once more
  \item \emph{Say it:} say \emph{pacās}
  \item \emph{Recall:} say the Gujarati for thirty, then the Gujarati for forty, then say \emph{pacās} again
\end{itemize}

\begin{tcolorbox}[breakable,colback=teal!4,colframe=teal!35!black,title={Before you move on}]
What does \gu{પચાસ} mean? (\textquotedblleft{}Fifty\textquotedblright{}.) Say it once more.
\end{tcolorbox}
\section[\texorpdfstring{\gu{સો} (so)}{સો (so)}]{\gu{સો} (so) — a hundred}

\label{lesson:GU-C83-hundred}



\begin{quote}
Before the new one: say the Gujarati for forty, then the Gujarati for fifty.
\end{quote}

\subsection*{You'll want to know: \gu{સો}}
\textbf{\gu{સો}} (\emph{so}) — \textquotedblleft{}a hundred\textquotedblright{}.

\begin{cousinweb}
\textbf{\gu{સો}} (\emph{so}) is a tadbhava, a word inherited from Sanskrit and reshaped by sound change. It comes from \emph{śata}, \textquotedblleft{}hundred\textquotedblright{} (CDIAL 12278).
\end{cousinweb}

\begin{grammarlens}[title={What you've built}]
One more word for fifty to a thousand, and how much.
\end{grammarlens}

\subsection*{Your turn}
\begin{itemize}
\raggedright
  \item \emph{Say it:} \emph{so}
  \item \emph{Say it:} \emph{so}, once more
  \item \emph{Say it:} say \emph{so}
  \item \emph{Recall:} say the Gujarati for forty, then the Gujarati for fifty, then say \emph{so} again
\end{itemize}

\begin{tcolorbox}[breakable,colback=teal!4,colframe=teal!35!black,title={Before you move on}]
What does \gu{સો} mean? (\textquotedblleft{}A hundred\textquotedblright{}.) Say it once more.
\end{tcolorbox}
\section[\texorpdfstring{\gu{હજાર} (hajār)}{હજાર (hajār)}]{\gu{હજાર} (hajār) — a thousand}

\label{lesson:GU-C83-thousand}



\begin{quote}
Before the new one: say the Gujarati for fifty, then the Gujarati for a hundred.
\end{quote}

\subsection*{You'll want to know: \gu{હજાર}}
\textbf{\gu{હજાર}} (\emph{hajār}) — \textquotedblleft{}a thousand\textquotedblright{}.

\begin{cousinweb}
\textbf{\gu{હજાર}} (\emph{hajār}) is a loan, usually traced to Persian \emph{hazār}, \textquotedblleft{}thousand\textquotedblright{}. Gujarati writes the \emph{z} as \emph{j}.
\end{cousinweb}

\begin{grammarlens}[title={What you've built}]
One more word for fifty to a thousand, and how much.
\end{grammarlens}

\subsection*{Your turn}
\begin{itemize}
\raggedright
  \item \emph{Say it:} \emph{hajār}
  \item \emph{Say it:} \emph{hajār}, once more
  \item \emph{Say it:} say \emph{hajār}
  \item \emph{Recall:} say the Gujarati for fifty, then the Gujarati for a hundred, then say \emph{hajār} again
\end{itemize}

\begin{tcolorbox}[breakable,colback=teal!4,colframe=teal!35!black,title={Before you move on}]
What does \gu{હજાર} mean? (\textquotedblleft{}A thousand\textquotedblright{}.) Say it once more.
\end{tcolorbox}
\section[\texorpdfstring{\gu{કેટલું} (keṭlũ)}{કેટલું (keṭlũ)}]{\gu{કેટલું} (keṭlũ) — how much}

\label{lesson:GU-C83-howmuch}



\begin{quote}
Before the new one: say the Gujarati for a hundred, then the Gujarati for a thousand.
\end{quote}

\subsection*{You'll want to know: \gu{કેટલું}}
\textbf{\gu{કેટલું}} (\emph{keṭlũ}) — \textquotedblleft{}how much\textquotedblright{}. The single form of \textbf{\gu{કેટલા}}.

\begin{cousinweb}
\textbf{\gu{કેટલું}} (\emph{keṭlũ}) is a tadbhava. CDIAL lists Gujarati \emph{keṭlī} and \emph{keṭlāk}, \textquotedblleft{}how much\textquotedblright{}, under a reconstructed \emph{*kiyatta}, \textquotedblleft{}how great\textquotedblright{} (CDIAL 3167).
\end{cousinweb}

\begin{grammarlens}[title={What you've built}]
One more word for fifty to a thousand, and how much.
\end{grammarlens}

\subsection*{Your turn}
\begin{itemize}
\raggedright
  \item \emph{Say it:} \emph{keṭlũ}
  \item \emph{Say it:} \emph{keṭlũ}, once more
  \item \emph{Say it:} say \emph{keṭlũ}
  \item \emph{Recall:} say the Gujarati for a hundred, then the Gujarati for a thousand, then say \emph{keṭlũ} again
\end{itemize}

\begin{tcolorbox}[breakable,colback=teal!4,colframe=teal!35!black,title={Before you move on}]
What does \gu{કેટલું} mean? (\textquotedblleft{}How much\textquotedblright{}.) Say it once more.
\end{tcolorbox}
\section[\texorpdfstring{\gu{કયું} (kayũ)}{કયું (kayũ)}]{\gu{કયું} (kayũ) — which one}

\label{lesson:GU-C83-which}



\begin{quote}
Before the new one: say the Gujarati for a thousand, then the Gujarati for how much.
\end{quote}

\subsection*{You'll want to know: \gu{કયું}}
\textbf{\gu{કયું}} (\emph{kayũ}) — \textquotedblleft{}which one\textquotedblright{}.

\begin{cousinweb}
\textbf{\gu{કયું}} (\emph{kayũ}) rests on the old question stem \emph{k-}. CDIAL has no entry for it, so the rest of its history is not settled.
\end{cousinweb}

\begin{grammarlens}[title={What you've built}]
One more word for fifty to a thousand, and how much.
\end{grammarlens}

\subsection*{Your turn}
\begin{itemize}
\raggedright
  \item \emph{Say it:} \emph{kayũ}
  \item \emph{Say it:} \emph{kayũ}, once more
  \item \emph{Say it:} say \emph{kayũ}
  \item \emph{Recall:} say the Gujarati for a thousand, then the Gujarati for how much, then say \emph{kayũ} again
\end{itemize}

\begin{tcolorbox}[breakable,colback=teal!4,colframe=teal!35!black,title={Before you move on}]
What does \gu{કયું} mean? (\textquotedblleft{}Which one\textquotedblright{}.) Say it once more.
\end{tcolorbox}
